\begin{problem}[第37届全国中学生物理竞赛决赛][反质子探测]{130}
    反粒子最早由狄拉克的理论所预言。1932年，安德森研究宇宙射线时发现了电子的反粒子——正电子。此后，人类又陆续发现了反质子等反粒子。
    用单粒子能量为$6.8 \mathrm{GeV}$的高能质子束轰击静止的质子靶，可产生反质子，其反应式为
    \begin{equation}
        p + p \rightarrow p + p + \underbrace {p + \overline {{p}}} _ {\longrightarrow} \pi^ {-} + \dots
        \label{eq:1.p130}
    \end{equation}
    质子和反质子湮灭时可产生$\pi ^ { - }$介子，因此在反质子束流中还伴随有大量$\pi ^ { - }$介子。图6a是探测反质子的实验装置原理图。反质子$\overline{\mathrm{p}}$和$\pi^{-}$介子流依次通过闪烁计数器$\mathrm{S}_{1}$、$\mathrm{S}_{2}$和$\mathrm{S}_{3}$，$\mathrm{S}_{1}$与$\mathrm{S}_{2}$相距$l = 12 \mathrm{m}$，在$\mathrm{S}_{2}$和$\mathrm{S}_{3}$之间放置有切伦科夫计数器$\mathrm{C}_{1}$和$\mathrm{C}_{2}$。切伦科夫计数器通过探测切伦科夫辐射（带电粒子在介质中的运动速度超过介质中光速时所激发的电磁辐射）而确定带电粒子的运动速度。$\mathrm{C}_{1}$仅记录速度较快的$\pi^{-}$介子，$\mathrm{C}_{2}$仅记录速度较慢的反质子。实验中$\mathrm{S}_{3}$的作用是检验前面的计数结果是否真实。
    \begin{subproblem}
        在上述反应中，假设末态质子和反质子速度近似相同，求反质子从$\mathrm{S}_{1}$运动到$\mathrm{S}_{2}$所需的时间$t _ { \bar { \mathrm { p } } }$。若$\pi ^ { - }$介子与反质子动能相同，求$\pi ^ { - }$介子从$\mathrm{S}_{1}$运动到$\mathrm{S}_{2}$所需的时间$t _ { \pi ^ { - } }$。
    \end{subproblem}
    \begin{subproblem}
        运动速率为$v \ ( v > c / n )$的带电粒子通过折射率为$n$的介质，求所产生的切伦科夫辐射传播方向与带电粒子的运动方向的夹角$\theta$。
    \end{subproblem}
    \begin{subproblem}
        微分式切伦科夫计数器可以记录速率在某一个区间的粒子数。图6b是其关于轴线（图中虚线）旋转对称的原理剖面图。光收集系统包括半径为$R$的球面镜和半径可调的环状光阑。当切伦科夫辐射传播方向与带电粒子运动方向的夹角$\theta$很小时，球面镜将带电粒子激发的切伦科夫辐射会聚在其焦平面上，形成半径为$r$的辐射光环，投射在计数器上。对于速率为$v$的带电粒子激发的辐射，求光环的半径。
    \end{subproblem}
    \begin{subproblem}
        反质子与质子相遇会发生湮灭反应$\mathrm{p} + \overline{\mathrm{p}} \to 3 \pi ^ { 0 }$。反应末态粒子的总动能与反应初态粒子的总动能之差即为反应能$Q$。为简单起见，假设质子和反质子的动能可忽略。末态三个$\pi ^ { 0 }$介子的总动能是一个常量，可用Dalitz图表示总动能在三个$\pi ^ { 0 }$介子之间的分配。如图6c所示，Dalitz图是一个高为1的等边三角形，$P$点到三边的距离等于三个$\pi ^ { 0 }$介子的动能占反应能的比率，即$d _ { i } = E _ { k , i } / Q$，$E _ { k , i }$表示第$i$个$\pi ^ { 0 }$介子的动能，$i =1,2,3$。以底边为$X$轴，底边中点为原点，底边上的中垂线为$Y$轴建立坐标系。求$P$点可能的分布范围边界的表达式，用反应能$Q$、$\pi ^ { 0 }$介子质量$m _ { \pi ^ { 0 } }$和真空中的光速$c$表示。并讨论若$m _ { \pi ^ { 0 } } = 0$时，$P$点的分布范围边界的表达式。
        已知：真空中的光速$c = 2.998 \times 10^{8} \mathrm{m/s}$，质子质量$M _ { \mathrm { p } } = 938.2720 \mathrm{MeV} / c ^ { 2 }$，中性$\pi$介子质量$m _ { \pi ^ { 0 } } = 134.9766 \mathrm{MeV} / c ^ { 2 }$，带电$\pi ^ { \pm }$介子质量$m _ { \pi ^ { \pm } } = 139.5702 \mathrm{MeV} / c ^ { 2 }$。
    \end{subproblem}
\end{problem}